% !TeX root = ../../main.tex
\subsubsection{Liang Cas13 fitness benchmark}

We used the transcriptome-scale RfxCas13d fitness screens of Liang
et al.\ across HAP1, HEK293FT, K562, MDA-MB-231, and THP1 cells
\citep{liang2026lncrna}.  The library contains 56,322 guides targeting 5,496
lncRNAs, 3,156 protein-coding genes, and 1,000 non-targeting controls.
The primary comparison used day 14 versus day 0 and the deposited normalized
guide values in Table S2, retaining two biological replicates where available.
Table S2C contains one K562 day-0 column and both day-14 columns.  We used this
deposited layout directly and applied no additional MAGeCK normalization.
These are fractional processed values after the authors' median-of-ratios
normalization, ComBat correction, and replicate-outlier procedure, not raw
integer read counts.  Because BARCS enforces integer counts, we explicitly
rounded each value to the nearest pseudo-count, recorded the rounding error,
and supplied the identical rounded matrix to BARCS, MAGeCK-RRA, and
MAGeCK-MLE.  Thus this analysis is explicitly a processed-count sensitivity
comparison: BARCS's beta-binomial and MAGeCK's negative-binomial likelihoods
do not retain their literal sampling interpretations.  Its purpose is to
compare rankings obtained from exactly the same public input, not to present
pseudo-count $p$-values as a raw-count validation.

For an optional raw-count confirmation, the repository also supplies a
restartable pipeline that streams endpoint reads from BioProject
PRJNA1344834 without retaining FASTQ files.  It implements the two stated
Cutadapt \citep{martin2011cutadapt} anchor rules and aligns 23-nucleotide
inserts to the Table S1B reference using Bowtie
\citep{langmead2009bowtie} arguments \texttt{-v 1 -m 3 --best -q}.  The article
prose describes ``up to three mismatches,'' whereas \texttt{-v 1} permits one;
the executable argument is therefore the reproducible contract.

Four gene-level analyses were kept distinct.  ``Liang RRA'' denotes the
authors' deposited day-14 results from \texttt{RobustRankAggreg} 1.2.1
\citep{kolde2012rra}; it is not MAGeCK-RRA.  Official MAGeCK-RRA used
\texttt{mageck test} without renormalizing Table S2, and official MAGeCK-MLE
used a day-14 coefficient with a
replicate indicator when two complete pairs were available
\citep{li2014mageck,li2015mageckvispr}.  BARCS used the identical
replicate-plus-day design, calibrated guide statistics with the 1,000
non-targeting controls, and applied the prespecified directional Stouffer
gene summary.  Each method retained its native gene statistic; agreement was
reported separately as Spearman correlation of gene effects.

The published Liang RRA calls were treated as a comparator, not ground truth.
The positive set comprised the library's known-essential protein-coding
genes, defined in the study from DepMap knockout screens.  Cell-line-specific
nulls were targeted lncRNAs with TPM equal to zero in both available
expression assays.  Primary operating metrics were average precision,
essential-control recall at an empirical 5\% null false-positive rate, the
absolute difference between the null $p<0.05$ fraction and its nominal 0.05
target, and directional essential-control recall at gene FDR 0.10.  Liang and
BARCS gene $p$-values were Benjamini--Hochberg adjusted within each cell line;
the official MAGeCK outputs supplied their gene FDR values.  These controls
test ranking and calibration without assuming that disagreement with the
published discovery list is an error.
